\documentclass[11pt,a4paper]{article}

\usepackage[a4paper,margin=23mm]{geometry}
\usepackage{iftex}
\ifPDFTeX
  \usepackage[T5]{fontenc}
  \usepackage[utf8]{inputenc}
  \usepackage[vietnamese]{babel}
  \usepackage{lmodern}
\else
  \usepackage{fontspec}
  \usepackage{polyglossia}
  \setmainlanguage{vietnamese}
  \setotherlanguage{english}
  \setmainfont{Times New Roman}
  \setsansfont{Arial}
  \setmonofont{Consolas}
\fi

\usepackage{xcolor}
\usepackage{microtype}
\usepackage{booktabs}
\usepackage{array}
\usepackage{longtable}
\usepackage{tabularx}
\usepackage{enumitem}
\usepackage{amsmath,amssymb}
\usepackage{graphicx}
\usepackage{tikz}
\usetikzlibrary{arrows.meta,positioning,fit,calc,shapes.geometric}
\usepackage{listings}
\usepackage{hyperref}
\usepackage{url}
\usepackage{fancyhdr}

\hypersetup{colorlinks=true,linkcolor=blue!55!black,urlcolor=blue!55!black}
\setlength{\emergencystretch}{3em}
\setlength{\parskip}{0.35em}
\setlength{\parindent}{1.2em}
\setlist{nosep,leftmargin=2em}
\renewcommand{\arraystretch}{1.16}

\definecolor{CoreBlue}{HTML}{155A9C}
\definecolor{CoreGreen}{HTML}{2E7D32}
\definecolor{CoreOrange}{HTML}{B85C00}
\definecolor{SoftBlue}{HTML}{EAF2FA}
\definecolor{SoftGray}{HTML}{F4F6F8}
\definecolor{LineGray}{HTML}{667085}

\lstdefinestyle{sourcecode}{
  basicstyle=\ttfamily\small,
  keywordstyle=\color{CoreBlue}\bfseries,
  commentstyle=\color{LineGray}\itshape,
  stringstyle=\color{CoreGreen},
  numbers=left,
  numberstyle=\tiny\color{LineGray},
  numbersep=7pt,
  frame=single,
  rulecolor=\color{LineGray!45},
  backgroundcolor=\color{SoftGray},
  breaklines=true,
  columns=fullflexible,
  keepspaces=true,
  showstringspaces=false,
  tabsize=4
}
\lstset{style=sourcecode}

\tikzset{
  flowbox/.style={rectangle,rounded corners=2pt,draw=LineGray,fill=SoftGray,
    align=center,text width=40mm,minimum height=13mm,font=\small,inner sep=5pt},
  flowwide/.style={rectangle,rounded corners=2pt,draw=LineGray,fill=SoftBlue,
    align=center,text width=53mm,minimum height=13mm,font=\small,inner sep=5pt},
  flowarrow/.style={-{Stealth[length=2.2mm]},thick,draw=CoreBlue},
  flowback/.style={{Stealth[length=2.2mm]}-{Stealth[length=2.2mm]},thick,draw=CoreGreen}
}

\pagestyle{fancy}
\fancyhf{}
\fancyhead[L]{Cẩm nang source VideoLevel}
\fancyhead[R]{Luồng native CAVLC và nhánh Python}
\fancyfoot[C]{\thepage}
\setlength{\headheight}{14pt}

\title{Cẩm nang đọc source và lần theo pipeline VideoLevel\\
\large Từ đầu vào H.264 đến proof, nhúng, blind extraction và xác minh}
\author{Tài liệu kỹ thuật tiếng Việt}
\date{26 tháng 9 năm 2026}

\begin{document}
\maketitle
\tableofcontents
\clearpage

\section{Mục đích và cách dùng tài liệu}

Tài liệu này là bản đồ đọc code của repository VideoLevel. Mục tiêu là giúp người
đọc lần theo một dữ liệu cụ thể qua từng lớp: script demo, giao thức API, chương
trình C++ phân tích CAVLC, bộ đóng gói Groth16 và bộ xác minh. Mỗi phần nêu file,
hàm, đầu vào, đầu ra, tác dụng lên pipeline và bước kế tiếp sử dụng kết quả đó
như thế nào. Đây là hướng dẫn source, không phải đặc tả cho mọi profile H.264.

Điểm quan trọng đầu tiên là repository hiện có hai đường xử lý video khác nhau.
Demo chính từ bước 00 đến 12 gọi CLI native C++ để nhúng/trích xuất theo lịch
blind có khóa. Script 13 chạy riêng nhánh Python cũ \path{src/embedder.py} /
\path{src/verifier.py}; nhánh này phân tích toàn video, dùng sidecar/vị trí và có
chi phí lớn hơn. Không nên lấy tên hàm Python \path{embed()} làm bằng chứng rằng
HTTP/WebSocket realtime hiện dùng hàm đó: native adapter gọi
\path{zkstego_blind_bits}.

\subsection{Ba lớp kết quả cần phân biệt}

\begin{itemize}
  \item \textbf{Thành công cú pháp video:} FFmpeg đọc và giải mã stream mà không
  báo lỗi nghiêm ngặt. Điều này chưa đo PSNR, SSIM hay chất lượng cảm nhận.
  \item \textbf{Thành công khôi phục payload:} extractor lấy đúng các byte đã
  đóng gói từ dấu hệ số CAVLC trong video. Điều này cần đúng khóa của lịch blind
  và frame authentication.
  \item \textbf{Thành công mệnh đề ZK:} verifier kiểm tra Groth16 proof trên đúng
  public inputs. Kết quả này không tự chứng minh nguồn gốc camera, danh tính
  người ký, thời hạn giấy phép hay video không bị xử lý ở nơi khác.
\end{itemize}

\subsection{Một khác biệt kiến trúc ảnh hưởng đến giao thức}

Video stego là artifact mang thông điệp và proof, nhưng trong cài đặt hiện tại
verifier còn cần khóa 32 byte để blind-extract và để helper Python dựng public
signals. Nếu verifier đã được cấp khóa an toàn từ trước thì media channel có thể
chỉ truyền video; nếu chưa, cần một kênh cấp khóa riêng. Endpoint native
\path{/api/v1/jobs/verify} hiện bị cấu hình thành 501: native API embed/extract
không tự chạy Groth16 verifier. Demo thực hiện Groth16 verification ở bước Python
riêng sau khi trích xuất.

\section{Sơ đồ hoạt động tổng thể}

\begin{figure}[ht]
\centering
\resizebox{\textwidth}{!}{%
\begin{tikzpicture}[node distance=9mm and 8mm]
  \node[flowbox] (inputs) {Video H.264 Annex-B\\message $m$\\khóa 32 byte $k$};
  \node[flowwide,right=of inputs] (proof) {\path{src/zk_proof.py}\\Groth16 proof cho payload\\proof nén còn 129 byte};
  \node[flowwide,right=of proof] (pack) {\path{pack()}\\$[\text{len}_{32BE}][m][proof]$\\ví dụ: 148 byte};
  \node[flowwide,right=of pack] (native) {CLI C++\\đọc IDR/CAVLC, chọn dấu theo $k$\\đóng frame xác thực và patch sign};
  \node[flowbox,right=of native] (video) {Video stego\\NAL H.264 đầu ra};
  \draw[flowarrow] (inputs) -- (proof);
  \draw[flowarrow] (proof) -- (pack);
  \draw[flowarrow] (pack) -- (native);
  \draw[flowarrow] (native) -- (video);

  \node[flowwide,below=23mm of native] (extract) {CLI C++ blind extractor\\video stego + khóa $k$\\đọc dấu theo cùng lịch, kiểm HMAC};
  \node[flowwide,left=of extract] (unpack) {\path{unpack()}\\tách message và proof 129 byte};
  \node[flowwide,left=of unpack] (verify) {\path{ZKSnarkBridge.verify()}\\dựng public signals\\Groth16 verify};
  \node[flowbox,left=of verify] (result) {Verifier\\ACCEPT / REJECT\\và message nếu hợp lệ};
  \draw[flowarrow] (video.south) -- ++(0,-8mm) -| (extract.north);
  \draw[flowarrow] (extract) -- node[above,font=\scriptsize]{payload byte[]} (unpack);
  \draw[flowarrow] (unpack) -- node[above,font=\scriptsize]{message, proof} (verify);
  \draw[flowarrow] (verify) -- (result);
  \draw[flowarrow,draw=CoreOrange] (inputs.south) -- ++(0,-18mm) -| node[pos=.25,below,font=\scriptsize,align=center]{khóa được cấp trước\\hoặc truyền ngoài video} (extract.south);
\end{tikzpicture}%
}
\caption{Luồng proof-bearing của demo. Khóa là input ngoài video ở cả pha nhúng lẫn blind extraction.}
\label{fig:overall}
\end{figure}

\subsection{Luồng gọi theo thứ tự}

\begin{longtable}{p{.07\textwidth}p{.24\textwidth}p{.29\textwidth}p{.28\textwidth}}
\toprule
\textbf{Bước} & \textbf{Điểm vào} & \textbf{Input chính} & \textbf{Output chuyển tiếp} \\
\midrule
\endfirsthead
\toprule\textbf{Bước} & \textbf{Điểm vào} & \textbf{Input chính} & \textbf{Output chuyển tiếp} \\
\midrule\endhead
00 & \path{demo/00_prepare.py} & Đường dẫn cover H.264, message text & \path{session.json}, \path{message.bin}, \path{secret_key.bin} \\
01 & \path{demo/01_generate_proof.py} & \path{message.bin}, \path{secret_key.bin}, artifacts circuit & proof JSON, proof nén 129 B, public signals \\
02 & \path{demo/02_pack_payload.py} & \path{message.bin}, proof nén & \path{packed_payload.bin} \\
03--04 & \path{demo/03_read_video.py}, \path{04_find_positions.py} & Annex-B H.264, kích thước payload & log parser, số IDR/candidate/capacity, probe một sign \\
05 & \path{demo/05_embed_video.py} & Cover, packed payload, key & \path{stego_video.h264} \\
06--07 & \path{demo/06_read_stego_video.py}, \path{07_blind_extract.py} & Stego, FFmpeg, key & log decode, \path{extracted_payload.bin} \\
08--09 & \path{demo/08_unpack_payload.py}, \path{09_verify_proof.py} & Payload, verification key, key & message/proof khôi phục và kết quả Groth16 \\
10--12 & Wrong-key, HTTP job, WebSocket E2E & Video/payload/key và channel test & Log từ chối, round-trip, test report \\
13 (tùy chọn) & \path{demo/13_python_pipeline.py} & Contract benchmark Python, video khác, key & python\_pipeline\_stego.h264, cache, log \\
14 & \path{demo/14_summary.py} & Các log hiện có và \path{session.json} & \path{SUMMARY.txt} \\
\bottomrule
\end{longtable}

\section{Cách đọc cấu trúc repository}

\begin{longtable}{p{.32\textwidth}p{.61\textwidth}}
\toprule
\textbf{Đường dẫn} & \textbf{Vai trò và ảnh hưởng đến pipeline} \\
\midrule
\endfirsthead
\toprule\textbf{Đường dẫn} & \textbf{Vai trò và ảnh hưởng đến pipeline} \\
\midrule\endhead
\path{demo/00...14_*.py} & Kịch bản có thể chạy theo từng công đoạn. Chúng ghép các module hiện có; không phải mỗi file là một implementation codec riêng. \\
\path{src/zk_proof.py} & Định nghĩa payload format, serialize proof Groth16, sinh witness/proof bằng Node/snarkjs, dựng public inputs và verify. \\
\path{circuits/payload_verify.circom} & Mệnh đề circuit: quan hệ SHA-256 giữa hash payload, commitment và secret; đây là logic được proving key chứng minh. \\
\path{native/tools/blind_bits.cpp} & CLI native: đọc key/payload, chọn chế độ, stream H.264, gọi API trong thư viện C++, ghi H.264 hoặc payload. \\
\path{native/src/cavlc_stream.cpp} & Parser Annex-B/SPS/PPS/IDR/CAVLC, sinh candidate sign, lịch keyed, framing HMAC, encoder/decoder incremental. \\
\path{native/include/zkstego/cavlc_stream.hpp} & Kiểu dữ liệu và interface công khai của C++: SPS/PPS/slice, macroblock, candidate, stream reader/encoder/decoder. \\
\path{src/api/native_handlers.py} & Adapter FastAPI/WebSocket gọi native CLI; quản lý input streaming, bounded buffers, resource metrics và HTTP handlers. \\
\path{src/api/app.py} & Lõi API chung: settings, auth token, giới hạn request/upload, queue/job store, file lưu tạm/cleanup, endpoint job. \\
\path{src/embedder.py}, \path{src/verifier.py} & Nhánh Python full-video cũ; phân tích cover, chọn/ghi sidecar vị trí, patch/reconstruct và verify theo cover. \\
\path{src/bitstream/} & Parser Python H.264/CAVLC và bit reader/writer; chủ yếu được dùng bởi nhánh Python, không phải parser native gọi từ FastAPI. \\
\path{src/core/pipeline.py}, \path{src/core/stego.py} & Phân tích IDR, vị trí safe, sửa/trích bit và dựng video cho nhánh Python; một số iterator streaming cũng được các prototype blind dùng. \\
\path{src/native_blind_contract.py} & Python executable specification của serialization candidate, HMAC schedule và frame native. Đây là contract/test aid; native CLI vẫn thi hành logic C++. \\
\path{src/realtime_cavlc.py} & Scheduler abstraction có giới hạn queue/latency. Tests exercise abstraction; native WebSocket adapter hiện tự quản lý subprocess/backpressure chứ không gọi scheduler này. \\
\path{src/key_policy.py}, \path{src/manifest.py} & Certificate thời hạn Ed25519 và sidecar manifest/signature. Đây không được gọi bởi demo native HTTP/WebSocket hiện tại; không nên kết luận nó tự hết hạn khóa native nếu chưa tích hợp tại enforcement point. \\
\path{src/trust/}, \path{src/provenance/} & Các helper fingerprint/attestation/provenance mở rộng. Không nằm trên luồng 00--12 đã demo; xem test/API cụ thể trước khi coi là runtime trust root. \\
\bottomrule
\end{longtable}

\section{Pha prover: từ message đến proof-bearing payload}

\subsection{Chuẩn bị session: script 00}

\path{demo/00_prepare.py::main()} parse cờ \path{--input} và \path{--message},
đọc toàn bộ cover để tính byte length/SHA-256, tạo folder
\path{demo/runs/session_<UTC>/}, ghi message UTF-8 thành \path{message.bin} và
sinh khóa demo ngẫu nhiên 32 byte vào \path{secret_key.bin} với quyền hạn chế.
\path{session.json} lưu đường dẫn input, số byte, hash và thời điểm; log không
ghi nội dung khóa.

Output này điều khiển mọi bước sau. Nếu đường dẫn \path{--session} không đúng,
\path{common.require_session()} dừng trước khi module proof/native được gọi. Việc
truyền \path{/runs/...} trên Windows tạo đường dẫn từ gốc ổ đĩa; từ thư mục
repository phải dùng đường dẫn đầy đủ hoặc \path{demo/runs/...}.

\subsection{Các helper dùng chung trong demo}

Mỗi file đánh số đều có \path{main()}, nhưng phần lớn chỉ điều phối: đọc
\path{--session}, kiểm artifact đầu vào, gọi module thật rồi in đường dẫn/log.
Các hành vi dùng chung nằm trong \path{demo/common.py}:

\begin{longtable}{p{.31\textwidth}p{.29\textwidth}p{.3\textwidth}}
\toprule\textbf{Hàm} & \textbf{Input $\to$ output} & \textbf{Vai trò trong stage} \\
\midrule\endfirsthead
\toprule\textbf{Hàm} & \textbf{Input $\to$ output} & \textbf{Vai trò trong stage} \\
\midrule\endhead
\path{require_session()} & Path argument $\to$ session folder + state dict & Resolve path, đọc \path{session.json}, dừng sớm nếu session hỏng/không tồn tại. \\
\path{read_key()} & Session $\to$ 32-byte key & Đọc key ở bước native/ZK; không ghi bytes khóa vào log. \\
\path{native_tool()} & Tên CLI $\to$ built executable path & Kiểm tra tìm thấy binary native đúng tên trước subprocess. \\
\path{require_artifact()} & Session + filename $\to$ existing path & Làm dependency gate giữa các bước, báo file trung gian còn thiếu. \\
\path{write_process_log()} & command/stdout/stderr/exit $\to$ text log & Lưu lệnh, mô tả stdin đã redacted, streams và mã thoát để debug. \\
\path{run_logged()} & Native command + optional stdin $\to$ completed process result & Chạy subprocess có timeout, lưu output/log; key stdin được mô tả redacted. \\
\path{python_stage_log()} & Stage/title/lines $\to$ stage log & Ghi marker kết quả cho script Python không có process con. \\
\path{run_python_logged()} & Callable Python $\to$ PASS/error log & Bắt exception, ghi traceback/status và cho stage dừng đúng mã lỗi. \\
\bottomrule\end{longtable}

\path{demo/test_common.py::test_stage_log_records_full_process_result_without_stdin_secret()}
kiểm tra log chứa đầy đủ kết quả process mà không lộ secret stdin. Các script
00--14 và artifact cụ thể của chúng được liệt kê ở Bảng 1; file test này chỉ
kiểm tra tiện ích log, không sinh artifact của một lần chạy pipeline.

\subsection{Circuit hiện đang chứng minh điều gì}

Trong \path{src/zk_proof.py}, payload là bytes $p$ và khóa 32 byte là $k$.
\path{ZKSnarkBridge._build_circuit_input()} tính
\[
 h=\operatorname{SHA256}(p), \qquad c=\operatorname{SHA256}(h\Vert k).
\]
Circuit \path{circuits/payload_verify.circom} nhận $h$, $c$ và độ dài payload làm
public signals, còn $k$ là private witness; constraints kiểm tra rằng $c$ đúng
là hash của $h\Vert k$ và độ dài nằm trong miền hợp lệ. Groth16 proof là bằng
chứng cho quan hệ này dưới circuit/proving key hiện hành. Circuit này không
nhận pixel/frame H.264 và không chứng minh video xuất phát từ camera hay rằng
video chưa bị transcode.

Trong demo, \path{01_generate_proof.py} gọi
\path{generate_proof_for_payload(message, key)} với đúng bytes trong
\path{message.bin}. Vì vậy proof ràng buộc message user (15 byte trong session
mẫu), không phải toàn bộ blob 148 byte đã gồm proof; bước 02 mới tạo outer blob
để chuyển qua H.264.

\begin{longtable}{p{.27\textwidth}p{.3\textwidth}p{.34\textwidth}}
\toprule
\textbf{Hàm/file} & \textbf{Input $\to$ output} & \textbf{Tác dụng kế tiếp} \\
\midrule
\endfirsthead
\toprule\textbf{Hàm/file} & \textbf{Input $\to$ output} & \textbf{Tác dụng kế tiếp} \\
\midrule\endhead
\path{_build_circuit_input(p,k)} & Bytes $p$, key 32 B $k$ $\to$ dictionary các bit $h,c,k$ và length & Node witness generator đọc JSON này. \\
\path{_compute_witness(input)} & Circuit input $\to$ file witness tạm & Gọi JS/WASM; xóa file input chứa key sau khi tạo witness. \\
\path{_snarkjs_prove(witness)} & Witness + proving key $\to$ proof JSON và public JSON & Trả về hai đối tượng snarkjs cho bước serialize và demo log. \\
\path{generate_proof_for_payload(p,k)} & Payload bytes + key $\to$ $(proof,public)$ & Đầu vào công khai ở đây không tự được gửi qua C++ embed; proof mới được serialize thành bytes. \\
\path{proof_to_bytes(proof)} & Pi A/B/C dạng point $\to$ 129 bytes & Cắt kích thước proof JSON để nhúng vào elementary stream. \\
\path{pack(p,proofBytes)} & $p$, proof 129 B $\to [\text{4-byte len}_{BE}][p][proof]$ & Blob được gửi cho native embed như payload opaque. \\
\path{blob_bit_length(p)} & Độ dài message $\to (4+|p|+129)\times8$ bit & Ước lượng capacity ở các caller Python cũ. \\
\bottomrule
\end{longtable}

\subsection{Ví dụ byte-for-byte với demo}

Với message \path{ZK demo payload}, UTF-8 dài 15 byte. Sau khi sinh proof,
\path{proof_to_bytes()} trả về đúng 129 byte. \path{pack()} ghép:
\begin{lstlisting}[numbers=none]
00 00 00 0f | 5a 4b 20 64 65 6d 6f 20 70 61 79 6c 6f 61 64 | <129 proof bytes>
  length=15                   message UTF-8
\end{lstlisting}
Tổng blob là $4+15+129=148$ byte, tương đương 1.184 bit. Blob 148 byte này
chính là đầu vào payload mà script 05 đưa cho native CAVLC embedder.

\subsection{Mã hóa/nén proof}

\path{proof_to_bytes()} ghi bốn tọa độ $x$ của các điểm Groth16 (mỗi tọa độ 32
byte big-endian) và một byte cờ chứa parity/sign của tọa độ $y$. Tổng là
$4\times32+1=129$ byte. \path{bytes_to_proof()} làm chiều ngược: khôi phục sign,
giải phương trình đường cong BN128 để dựng lại $y$, rồi tạo cấu trúc pi\_a,
pi\_b, pi\_c mà snarkjs nhận.

Đây là serialization riêng của source. Không được cắt/gắn proof tùy ý: mọi byte
đều đi vào phép giải nén tọa độ hoặc cờ sign, và sai định dạng làm verify thất
bại hoặc bị assertion/error.

\section{Đường native: đọc H.264 và chọn vị trí sign CAVLC}

\subsection{Tại sao sửa đúng một bit hệ số có thể giữ parser đồng bộ}

CAVLC residual block mở đầu bằng coeff\_token để mã hóa TotalCoeff và
TrailingOnes. Với mỗi trailing one, stream ghi một sign bit riêng: bit 0 nghĩa
giá trị hệ số dương 1, bit 1 nghĩa âm 1. Đổi sign chỉ thay giá trị hệ số từ $+1$
sang $-1$ hoặc ngược lại; không đổi số hệ số, số zero-run, VLC codeword length
hay vị trí bắt đầu block tiếp theo. Vì vậy native implementation tìm sign
position đã parse được rồi patch đúng RBSP bit đó, thay vì decode frame thành
pixel rồi encode lại.

Điều này chỉ an toàn cho candidate sign hợp lệ. Không thể lật tùy tiện một bit
trong NAL: nó có thể sửa coeff\_token, level, macroblock header hoặc emulation
prevention byte và khiến các bit sau bị đọc lệch.

\subsection{Annex-B, EBSP và RBSP}

\path{split_annex_b()} quét start code 3/4 byte, cắt elementary stream thành NAL
units và giữ header/index/offset. \path{AnnexBNalUnit::rbsp()} gọi
\path{ebsp_to_rbsp()} để bỏ emulation-prevention byte \path{0x03} đặt sau
\path{00 00}. Parser đọc syntax trên RBSP. Khi đã patch xong,
\path{rbsp_to_ebsp()} chèn prevention bytes cần thiết và
\path{assemble_annex_b()} ghép start code + header + payload về Annex-B.

Ví dụ nếu trong EBSP có chuỗi byte mã hóa tương đương \path{00 00 01}, encoder
phải ngăn chuỗi đó bị nhầm với start code; parser trước hết bỏ byte escape, đọc
bit syntax trên RBSP, sau đó serializer phục hồi escape khi ghi output.

\subsection{Các kiểu dữ liệu C++ giữ kết quả parser}

\begin{longtable}{p{.32\textwidth}p{.58\textwidth}}
\toprule
\textbf{Kiểu trong header} & \textbf{Dữ liệu và consumer tiếp theo} \\
\midrule
\endfirsthead
\toprule\textbf{Kiểu trong header} & \textbf{Dữ liệu và consumer tiếp theo} \\
\midrule\endhead
\path{AnnexBNalUnit} & Offset, start-code length, NAL header/type, EBSP payload; dùng để tìm IDR, parameter set và RBSP. \\
\path{H264BaselineSps} & Profile/level, kích thước picture theo macroblock, frame/POC fields; xác định cấu trúc và cách đọc header. \\
\path{H264BaselinePps} & Entropy mode và các syntax flag; native path yêu cầu CAVLC, không phải CABAC. \\
\shortstack[l]{\texttt{H264BaselineIdr}\\\texttt{SliceHeader}} & first MB, slice type, PPS id, frame/id fields, QP delta và offset residual đầu slice. \\
\path{CavlcCoeffToken} & TotalCoeff, TrailingOnes và offset của sign bits; là nền để giải mã level/tail. \\
\path{CavlcDecodedLumaBlock} & Token, levels, zero/run tail, hệ số tái dựng; consumer tạo sign candidates. \\
\path{CavlcDecodedIMacroblock} & Header I-macroblock và residual plane blocks; slice decoder đi qua macroblock theo địa chỉ. \\
\path{CavlcDecodedIdrSlice} & SPS/PPS/slice header, macroblocks và trailing-bit offset; đầu vào collector candidate. \\
\path{CavlcSignCandidate} & (NAL index, macroblock address, category, block index, RBSP bit offset); identity vật lý sẽ được xếp thứ tự bằng key. \\
\shortstack[l]{\texttt{AuthenticatedCavlc}\\\texttt{StreamSegment}} & Bytes output, candidate capacity, bit count đã nhúng, trạng thái complete; trả về qua mỗi segment. \\
\bottomrule
\end{longtable}

\subsection{Call chain C++ từ Annex-B đến candidate}

Các hàm dưới đây nằm chủ yếu trong \path{native/src/cavlc_stream.cpp}, interface
ở \path{native/include/zkstego/cavlc_stream.hpp}:

\begin{longtable}{p{.32\textwidth}p{.29\textwidth}p{.3\textwidth}}
\toprule
\textbf{Hàm} & \textbf{Input $\to$ output} & \textbf{Ảnh hưởng/consumer} \\
\midrule
\endfirsthead
\toprule\textbf{Hàm} & \textbf{Input $\to$ output} & \textbf{Ảnh hưởng/consumer} \\
\midrule\endhead
\path{ebsp_to_rbsp()} / \path{rbsp_to_ebsp()} & Byte vector $\leftrightarrow$ byte vector & Chuyển giữa byte syntax NAL bên ngoài và bitstream RBSP parser/patcher thao tác. \\
\path{parse_baseline_sps()} & SPS RBSP $\to$ SPS struct & Cung cấp dimensions/flags cho slice và MB parsing. \\
\path{parse_baseline_pps()} & PPS RBSP $\to$ PPS struct & Cung cấp PPS syntax; kiểm entropy mode và parameters. \\
\path{parse_baseline_idr_slice_header()} & Slice RBSP + SPS/PPS $\to$ header, data offset & Chỉ ra điểm bắt đầu macroblock residual. \\
\path{parse_cavlc_coeff_token()} & RBSP, start bit, $n_C$ $\to$ TotalCoeff/TrailingOnes/sign offsets & Chọn bảng VLC đúng context; trả candidate sign offsets. \\
\path{decode_cavlc_non_trailing_levels()} & Token + RBSP $\to$ trailing values, levels, next offset & Nhảy qua phần level sau sign bits để xác định tail đúng. \\
\path{decode_cavlc_residual_tail()} & TotalCoeff, giới hạn block, offset $\to$ total zeros/runs/next offset & Đồng bộ vị trí block kế tiếp. \\
\path{reconstruct_cavlc_block()} & Levels + runs + max coefficients $\to$ vector coefficients & Xác nhận residual có cấu trúc phù hợp và điền vị trí coefficient. \\
\path{decode_cavlc_luma_block()} / \path{decode_cavlc_chroma_*()} & RBSP + offset + context $\to$ decoded block & Chuyên biệt hóa luma/chroma, luma AC/DC. \\
\path{decode_cavlc_luma_macroblock()} & CBP, offset, neighbour TotalCoeff $\to$ 16 luma blocks & Tính $n_C$ theo hàng xóm đã giải mã, xử lý từng block coded. \\
\path{parse_baseline_i_macroblock_header()} & RBSP + offset $\to$ MB type, CBP, QP delta, residual offset & MB parser quyết định syntax residual tiếp theo. \\
\path{decode_baseline_i_idr_slice()} & RBSP + SPS/PPS $\to$ một decoded IDR slice & Lặp macroblock I slice và gom offset/sign metadata. \\
\path{decode_baseline_i_idr_slices()} & Annex-B bytes $\to$ danh sách IDR I slices & Dùng bởi capacity scan, embedding và extraction full-buffer. \\
\path{collect_cavlc_trailing_one_sign_candidates()} & Decoded slices $\to$ candidate vector & Mỗi residual block đóng góp các sign offsets trailing-one patchable theo parser. \\
\path{inspect_baseline_idr_headers()} & Annex-B $\to$ header/mẫu đầu IDR & Công cụ inspect dùng để giải thích NAL/MB đầu, không quyết định lịch embed cuối. \\
\bottomrule
\end{longtable}

\subsection{Ví dụ một candidate}

Candidate trong log inspect demo có thể đọc như địa chỉ trong syntax, không phải
tọa độ pixel: \path{patched_nal=3 macroblock=0 rbsp_bit=91}. Công cụ
\path{zkstego_idr_inspect --flip-first-sign} tạo \path{candidate_probe.h264}
để thấy một sign bit riêng lẻ có thể patch. File probe này không chứa payload
hoàn chỉnh và không phải output của bước 05.

\section{Lịch keyed, frame integrity và stream encoder/decoder}

\subsection{Định danh candidate và thứ tự bí mật}

\path{serialize_cavlc_sign_candidate()} biến tuple candidate thành chuỗi định
danh ổn định gồm NAL, MB, loại residual, block và offset. Hàm
\path{score_keyed_cavlc_sign_candidate()} tính HMAC-SHA-256 trên định danh bằng
khóa lịch có domain separation. \path{select_keyed_cavlc_sign_candidates()} sắp
xếp candidate theo score rồi lấy đủ số bit cần. Verifier quét cùng video stego,
tái tạo cùng candidate list và dùng cùng key, nên trích bit theo đúng thứ tự.
Khóa không được giấu trong H.264; nó phải được provision riêng.

\subsection{Authenticated frame bên trong stream}

Native C++ dùng một framing khác với outer payload format của Python:
\begin{lstlisting}[numbers=none]
frame = version_u8 || payload_length_u16_be || payload || tag_16_bytes
\end{lstlisting}
Tag là 16 byte đầu của HMAC-SHA-256 theo domain riêng. Phần overhead là
$1+2+16=19$ byte. Vì vậy blob proof-bearing 148 byte trở thành frame 167 byte,
tức 1.336 bit phải nhúng. Đây là lý do log capacity báo 1.336 bit, trong khi
\path{02_pack_payload.log} báo 1.184 bit.

\path{pack_authenticated_cavlc_frame()} thêm version/length/tag;
\path{unpack_authenticated_cavlc_frame()} kiểm version, giới hạn length, đúng
kích thước và so sánh tag constant-time trước khi trả payload. HMAC không thay
Groth16: HMAC bảo vệ framing/channel với shared key; Groth16 kiểm quan hệ
payload/key theo circuit. Wrong key có thể hỏng keyed schedule và/hoặc HMAC
frame, trước cả khi đến bước Groth16.

\subsection{Patch và extraction bit-level}

\path{embed_keyed_cavlc_sign_bits()} nhận Annex-B, key, bit vector; giải mã IDR,
tạo candidate, chọn đúng số sign offsets và patch sign bit trong RBSP. Hàm không
re-encode nguyên video. \path{extract_keyed_cavlc_sign_bits()} tái tạo candidate
list và đọc dấu theo cùng thứ tự. Các helper
\path{bytes_to_msb_bits()} / \path{msb_bits_to_bytes()} quy định bit order MSB
trước để hai phía đồng bộ.

Ở chế độ stream, \path{AuthenticatedCavlcStreamEncoder::process_segment()} nhận
một Annex-B segment (SPS/PPS cộng IDR nếu có), lấy tối đa cấu hình bit mỗi IDR,
patch phần kế tiếp của framed payload và báo \path{bits_embedded} cùng trạng
thái complete. \path{complete()} / \path{remaining_bits()} hỗ trợ caller biết
khi nào payload xong. Decoder dùng
\path{AuthenticatedCavlcStreamDecoder::consume_segment()} để tích lũy sign bits,
đọc header length, chờ đủ body+tag rồi chỉ xuất
\path{authenticated_payload()} sau khi kiểm tag; \path{failed()} là trạng thái
fail-closed.

\subsection{CLI native và chế độ byte stream}

File \path{native/tools/blind_bits.cpp} là process boundary mà Python API gọi.
\path{main()} dispatch lệnh như \path{measure-live-capacity-stdin},
\path{embed-stream-auth-stdin}, \path{extract-stream-auth} và
\path{extract-live-auth-stdin}. Các helper chính:

\begin{longtable}{p{.34\textwidth}p{.26\textwidth}p{.29\textwidth}}
\toprule\textbf{Hàm} & \textbf{I/O} & \textbf{Vai trò} \\
\midrule\endfirsthead
\toprule\textbf{Hàm} & \textbf{I/O} & \textbf{Vai trò} \\
\midrule\endhead
\path{decode_hex()} & Hex key/payload $\to$ sensitive bytes & Kiểm tra ký tự hex và tránh in secret. \\
\path{read_key_file()} / \path{read_key_stdin_line()} & Key source $\to$ key bytes & Hỗ trợ file quyền hạn chế hoặc truyền key qua stdin, không argv. \\
\path{read_binary()} / \path{write_new_binary()} & File bytes $\leftrightarrow$ vector & I/O command non-live; output tạo mới để không ghi đè. \\
\path{run_live_capacity_measure()} & Annex-B stdin $\to$ JSON metrics stdout & Parse stream và đếm IDR/candidate/capacity. \\
\path{run_live_authenticated_embed()} & stdin chứa key, payload rồi Annex-B $\to$ Annex-B stego stdout & Incremental reader + stream encoder; xử lý từng NAL/IDR. \\
\path{run_live_authenticated_extract()} & Key line + Annex-B stdin $\to$ payload hex stdout & Incremental reader + stream decoder; kết thúc khi payload authenticate. \\
\path{set_standard_streams_to_binary()} & stdin/stdout handles $\to$ binary mode & Bắt buộc trên Windows để CRLF/encoding text không đổi bytes H.264. \\
\path{parse_bit_count()} & decimal string $\to$ \path{size_t} & Reject số âm, overflow, ký tự thừa. \\
\bottomrule\end{longtable}

\path{AnnexBNalStreamReader::read_next()} giữ buffer hữu hạn quanh một NAL và một
input chunk, nhận cả start code đi qua ranh giới chunk; giới hạn NAL tối đa mặc
định 16 MiB. Đây là cơ chế giảm yêu cầu giữ toàn bộ raw input trong live
WebSocket, nhưng không đồng nghĩa toàn bộ API/circuit/video analytics đều có
memory bounded như nhau.

\section{Verifier: từ video stego đến quyết định accept/reject}

\subsection{Blind extraction native}

\path{demo/07_blind_extract.py::main()} đọc \path{stego_video.h264}, lấy key
từ session, gọi \path{zkstego_blind_bits extract-stream-auth}. Cover video không
được truyền cho lệnh này. Output stdout hex được chuyển về bytes và ghi
\path{extracted_payload.bin}. Native decoder:

\begin{enumerate}
  \item cập nhật SPS/PPS gần nhất;
  \item nhận IDR NAL, dựng segment cho stream decoder;
  \item decode lại CAVLC candidates, xếp theo keyed schedule, gom sign bits;
  \item đọc frame length/version, đợi đủ dữ liệu, xác thực HMAC;
  \item trả payload bytes hoặc lỗi nếu stream hết/sai key/sai tag.
\end{enumerate}

\subsection{Unpack và verify Groth16}

\path{unpack(blob)} kiểm tra tối thiểu 4 byte đầu, đọc độ dài message big-endian,
kiểm tra blob đủ $4+|m|+129$ byte rồi tách $(m,proofBytes)$. Trong demo,
\path{08_unpack_payload.py} ghi hai output riêng
\path{recovered_message.bin} và \path{recovered_proof_129_bytes.bin}.

Tiếp theo \path{bytes_to_proof()} tái dựng điểm Groth16.
\path{ZKSnarkBridge._build_public_signals(message,key)} tính hash message,
commitment và length theo circuit. \path{verify()} gọi snarkjs với
\path{verification_key.json}; \path{_snarkjs_verify()} tạo temporary public/proof
JSON rồi đọc exit/result. \path{verify_proof_for_payload()} là wrapper thuận
tiện dựng public signals từ payload/key trước khi gọi verify.

Trong demo 09, cùng proof với message gốc phải valid; nối thêm một byte
\path{!} phải invalid. Đây là test negative hữu ích, nhưng không kiểm tra mọi
kiểu tamper vào H.264, mọi public statement hoặc quản lý key trong production.

\subsection{HTTP và WebSocket làm gì, không làm gì}

\path{create_native_app()} trong \path{src/api/native_handlers.py} dựng app bằng
\path{create_app()} nhưng đưa \path{verify_handler=None}. HTTP embed nhận cover,
message bytes và secret key; HTTP extract nhận stego cùng key. Endpoint verify
native vì thế trả 501. Nếu muốn embed Groth16, caller phải tự tạo proof-bearing
blob trước và truyền blob đó như \path{message}; native layer coi payload là
opaque bytes.

WebSocket \path{/api/v1/stream} nhận control JSON đầu tiên. Embed control chứa
operation, key và message base64; sau đó client đẩy các binary chunk Annex-B và
gửi control \path{"type":"end"}. Server mở native process, forward chunk qua
stdin và output qua WebSocket với flow-control/backpressure. Extract control
chứa operation, key và maximum payload bytes; output là extracted bytes.
\path{native_stream()} giới hạn thời gian, tổng byte, kích thước chunk, số chunk,
stdin high/low-water mark và lấy mẫu RSS/CPU/latency. Đây là service streaming
trên elementary H.264; input camera device/capture driver chưa tự được mở bởi
module này.

\section{API chung: hàm nào nhận request, hàm nào chạy job}

\path{src/api/app.py::create_app()} đăng ký endpoints và ghép các handler được
inject. Các khối hàm chính:

\begin{longtable}{p{.34\textwidth}p{.28\textwidth}p{.27\textwidth}}
\toprule\textbf{Hàm/lớp} & \textbf{Input $\to$ output} & \textbf{Ảnh hưởng} \\
\midrule\endfirsthead
\toprule\textbf{Hàm/lớp} & \textbf{Input $\to$ output} & \textbf{Ảnh hưởng} \\
\midrule\endhead
\path{RequestBodyGuard.__call__()} & ASGI receive/send $\to$ guarded request & Reject body quá giới hạn hoặc timeout trước khi parse upload. \\
\path{WorkDirLease.acquire/release()} & work directory $\to$ lease lifecycle & Ngăn nhiều service instance dùng chung work directory trái phép. \\
\path{ApiSettings.from_environment()} & environment $\to$ settings & Nạp token, paths, queue/upload/time limits. \\
\path{JobStore.create/get/update/discard()} & job record $\leftrightarrow$ JSON on disk & Job ID/status persist để client poll và download. \\
\path{_decode_key()} / \path{_decode_message()} & base64 form field $\to$ bytes & Kiểm format/độ dài trước gọi native handler. \\
\path{_save_upload()} & UploadFile $\to$ file bounded-size & Ghi video input theo giới hạn request; lỗi sẽ cleanup. \\
\path{require_token()} & Authorization header $\to$ pass/401 & API bearer authentication, không phải key CAVLC/Groth16. \\
\path{reserve_capacity()} & semaphore $\to$ slot hoặc reject & Chặn vượt số worker/job queue cấu hình. \\
\path{submit()} & job id + operation + callable $\to$ worker thread, cập nhật state & Chạy tác vụ nền; job endpoint trả 202 rồi client poll. \\
\path{startup()/shutdown()} & service lifecycle $\to$ recovery/cleanup & Khôi phục job dang dở, khởi động/dừng cleanup worker. \\
\path{start_embed()} & upload cover + message/key $\to$ accepted job & Lưu input, tạo job và gọi injected embed handler. \\
\path{start_extract()} & upload stego + key/bounds $\to$ accepted job & Gọi extractor và lưu payload artifact để download. \\
\path{start_verify()} & stego + cover + key/length $\to$ job hoặc 501 & Chỉ chạy nếu verify handler được inject; native adapter cố ý không inject. \\
\path{get_job()} / \path{download_artifact()} & job id $\to$ status/file response & Poll và tải kết quả; cleanup/policy quyết định thời gian tồn tại. \\
\bottomrule\end{longtable}

Trong \path{native_handlers.py}, \path{invoke()} đưa key vào stdin child process,
không đặt key trên command line hay key temp file. \path{embed_handler()} gọi
\path{embed-stream-auth-stdin}; \path{extract_handler()} gọi
\path{extract-stream-auth} rồi validate hex, giới hạn payload và tạo file mới.
\path{native_stream()} là đường WebSocket độc lập với queue HTTP; nó quản lý
child process, tasks cho stdin/stdout, cancellation, deadlines và metrics.

\section{Nhánh Python full-video cũ: các hàm và phụ thuộc}

\subsection{Không nhầm nhánh này với native blind}

\path{src/embedder.py::embed()} là pipeline Python đã tồn tại song song. Nó cần
cover để phân tích vị trí, tạo proof bên trong hàm, gói payload, lấy safe
positions, patch coefficients qua Python reconstruction và ghi sidecars. Hàm
\path{src/verifier.py::verify()} yêu cầu cả stego lẫn original cover, trừ khi
caller cung cấp vị trí đã tính trước. Demo 13 gọi nhánh này riêng trên operating
contract benchmark khác.

\subsection{Các helper của embedder}

\begin{longtable}{p{.34\textwidth}p{.28\textwidth}p{.27\textwidth}}
\toprule\textbf{Hàm trong \path{src/embedder.py}} & \textbf{Input $\to$ output} & \textbf{Consumer} \\
\midrule\endfirsthead
\toprule\textbf{Hàm} & \textbf{Input $\to$ output} & \textbf{Consumer} \\
\midrule\endhead
\path{_get_bridge(circuits_dir)} & path $\to$ cached bridge & Dùng lại ZK wrapper cho lần embed tiếp theo. \\
\path{_prune_patchable_positions()} & safe positions + offsets $\to$ chỉ candidate block có thể patch & Ngăn yêu cầu vượt quá patchability thật của bitstream. \\
\path{_limit_positions_per_block()} & positions, cap $\to$ bounded schedule & Thực thi giới hạn sửa mỗi block để giảm artifact/duplicate. \\
\path{_filter_reconstructed_positions()} & planned/applied block sets $\to$ thật sự được output áp dụng & Không báo “đã nhúng” cho patch mà reconstructor bỏ qua. \\
\path{_resolve_video_analysis()} & video hoặc precomputed tuple $\to$ 6 phần phân tích & Ưu tiên tuple caller, nếu không nạp/build cache. \\
\path{embed()} & video, message, key, circuit path, options $\to$ EmbedResult + H.264 + sidecars & Orchestrator của Python legacy branch. \\
\bottomrule\end{longtable}

Bên trong \path{embed()}, trình tự chính là: validate file/message/key;
\path{analyze_stream_profile()}; sinh Groth16 proof; nén proof; \path{pack()};
(tùy chọn) chaos scramble; resolve analysis/cache; lọc safety/patchability; gọi
\path{PayloadEmbedder.embed_payload()}; gọi
\path{BitstreamReconstructor.reconstruct_video()}; lọc lại vị trí theo
\path{applied_block_keys}; thử tối đa 12 vòng loại các block không tái dựng
được; lưu video, \path{.positions.json}, \path{.meta.json},
\path{.manifest.json}; trả \path{EmbedResult}. Dữ liệu quan trọng trong result
gồm bits yêu cầu/đã nhúng, capacity từng lớp, output path, proof/public dict và
used positions.

\subsection{Phân tích H.264/CAVLC của Python path}

\begin{longtable}{p{.35\textwidth}p{.26\textwidth}p{.28\textwidth}}
\toprule\textbf{Hàm/lớp} & \textbf{Input $\to$ output} & \textbf{Ý nghĩa trong chuỗi} \\
\midrule\endfirsthead
\toprule\textbf{Hàm/lớp} & \textbf{Input $\to$ output} & \textbf{Ý nghĩa trong chuỗi} \\
\midrule\endhead
\path{NALParser.parse()} & H.264 bytes $\to$ NALUnit list & Tách start codes, header và EBSP để các parser cấp sau làm việc. \\
\path{_find_start_codes()} & bytes $\to$ [(offset,size)] & Tìm phân cách Annex-B 3/4 byte. \\
\path{_extract_nal_unit()} & offset range $\to$ NALUnit & Lấy type/ref idc và loại emulation-prevention bytes cho RBSP. \\
\path{SliceHeaderParser.parse()} & reader + NAL/SPS/PPS $\to$ SliceHeader & Đọc các Exp-Golomb/slice flags cần để đến macroblock data. \\
\path{MacroblockParser.parse_macroblock()} & reader + slice context $\to$ MacroblockData & Đọc mb type, prediction/CBP và phân ranh residual syntax. \\
\path{MacroblockParser.calculate_nC()} & tọa độ block + neighbour counts $\to n_C$ & Context này chọn đúng coeff\_token table của CAVLC. \\
\path{TraceableCAVLCParser.extract_with_offsets()} & IDR NAL + SPS/PPS $\to$ block levels, offsets, MB indices & Cấp cho safety filter/patcher biết bit nào mã hóa hệ số nào. \\
\path{CAVLCDecoder.decode_block_cavlc()} & bit reader + $n_C$ + max coeff $\to$ CoefficientBlock & Giải coeff\_token, trailing signs, levels và zero runs thành coefficient list. \\
\path{_decode_coeff_token()} & $n_C$ + reader $\to$ TotalCoeff, TrailingOnes & TrailingOnes quyết định sign bits được đọc trước levels. \\
\path{_decode_levels()} & count + trailing ones $\to$ level list & Đọc level code sau sign prefix. \\
\path{_decode_total_zeros()} / \path{_decode_runs()} & block residual syntax $\to$ positions of nonzero coeff & Giữ đúng scan order khi reconstruct block. \\
\path{CAVLCEncoder.encode_block_cavlc()} & coefficient list + $n_C$ $\to$ VLC bits & Legacy reconstruction reserializes modified blocks. \\
\path{BitstreamReader.read_ue/read_se()} & RBSP bits $\to$ unsigned/signed Exp-Golomb value & Dùng cho syntax H.264 như ue(v), se(v). \\
\bottomrule\end{longtable}

\subsection{Safety filter, patch và tái dựng Python}

Trong \path{src/core/stego.py}, \path{CAVLCSafetyFilter._detect_trailing_ones()}
tìm index trailing-one; \path{_verify_block_bit_length_invariance()} kiểm rằng
ứng viên không làm đổi bit length theo tiêu chí patch; \path{get_safe_positions()}
lấy tập an toàn rộng hơn; \path{get_safe_sign_positions()} giới hạn vào sign
positions phù hợp. \path{sort_blocks_interleaved()} phân bố thứ tự theo block/MB
để tránh dồn sửa một vùng.

\path{PayloadEmbedder.embed_payload()} chuyển blob thành bit, cấp từng bit cho
candidate và gọi \path{_embed_with_safety_filter()}; output là danh sách
coefficient block đã đổi, số bit thực nhúng và vị trí đã dùng. Chiều đọc là
\path{extract_payload()} / \path{_extract_with_safety_filter()}, lấy sign từ
coefficient theo đúng position order rồi \path{_bits_to_bytes()}. Hàm
\path{_modify_lsb()} thực chất sửa dấu/LSB theo quy tắc block hiện hành; không
nên hiểu là sửa pixel LSB trong ảnh RGB.

\path{BitstreamPatcher.patch_slice()} áp các modification vào NAL theo offset;
\path{validate_block_patchability()} xác nhận vị trí/độ dài block;
\path{_encode_coefficients_to_bits()} tái mã hóa CAVLC block. Sau đó
\path{BitstreamReconstructor.reconstruct_video()} cập nhật NALs và gọi
\path{_reconstruct_slice_with_cavlc()} / \path{_write_h264_file()} để tạo file
mới; SPS/PPS parser duy trì context; emulation prevention được thêm lại. Các
validator FFmpeg/PSNR là tùy chọn, không đồng nghĩa mọi lần chạy embed mặc định
đều đã chứng minh quality.

\subsection{Verifier Python yêu cầu những gì}

\path{src/verifier.py::verify()} validate cả \path{stego_video_path} và
\path{original_video_path}, circuit path, 32-byte key và message length. Nó đọc
sidecar qua \path{_load_sidecar_data()} nếu có; nếu không thì tái phân tích cover
để tạo position map. Tiếp đó tính số bit cần, dựng CAVLC analysis, áp chaos
inverse nếu bật, đọc bit bằng \path{extract_bits_direct()}, gọi \path{unpack()},
\path{bytes_to_proof()}, dựng public signals rồi \path{bridge.verify()}. Output
\path{VerifyResult} gồm valid, message (chỉ giữ khi hợp lệ), proof/public dict và
bits extracted.

\path{verifier_blind_keyed.py::verify_blind_keyed()} là một thử nghiệm riêng
không sidecar: lấy message length để ước lượng fixed payload bits, gọi
\path{derive_blind_positions_validated_pool_proxy()}, parse stego, trích block,
đọc bits, unpack và verify. Nó nêu assumption cố định/no-chaos; không nên đánh
đồng nhất với native stream decoder vốn có format/frame/selector contract khác.

\section{Cache, scheduler và những nhánh hỗ trợ}

\subsection{Cache phân tích video}

\path{src/core/analysis_cache.py::_video_fingerprint()} tạo dấu vân tay đầu vào;
\path{_code_fingerprint()} gắn cache với fingerprint source; \path{_cache_path()}
chọn file; \path{load_or_build_video_analysis()} chạy NAL/CAVLC analysis và
safety filter hoặc nạp trusted pickle; \path{_build_reconstruction_context()}
chuẩn bị dữ liệu reconstructor; \path{load_or_build_reconstruction_context()}
đọc/build context riêng; \path{clear_video_analysis_cache()} xóa các pickle cache.

Cache giảm thời gian lặp lại nhưng pickle chỉ an toàn nếu nguồn cache đáng tin
cậy. Biến \path{ZK_STEGO_TRUSTED_PICKLE_CACHE=1} là opt-in rõ ràng ở nhánh
Python demo; không bật với cache nhận từ nguồn lạ. Native CLI demo không dùng
Python pickle analysis cache.

\subsection{Realtime scheduler abstraction}

\path{CAVLCRealtimeBudget.__post_init__()} validate FPS/queue/memory bounds;
\path{frame_interval_s} trả $1/fps$; \path{latency_limit_s} chọn p95 limit hoặc
frame interval. \path{RealtimeCAVLCScheduler.submit()} validate segment size,
epoch/cache và queue bound; \path{process_next()} lấy segment + proof callback,
gọi patch callback, đo latency; \path{report()} tính processed/dropped/p95 và
\path{accepted}. Đây là orchestration/test abstraction; implementation comment
chủ động cấm dùng full-video Python embed làm patcher realtime. Native WebSocket
hiện gọi native process trực tiếp và có cơ chế riêng.

\subsection{Key certificate và manifest}

\path{src/key_policy.py::issue_key_certificate()} ký Ed25519 certificate có
subject public key, key id, issuer, not-before và expires-at;
\path{KeyCertificate.verify()} kiểm chữ ký issuer, format/version và cửa sổ
thời gian; \path{verify_manifest_with_certificate()} kiểm certificate trước,
rồi dùng subject public key để verify manifest.
\path{src/manifest.py::StegoManifest.sign()}, \path{verify_signature()},
\path{compute_content_hash()}, \path{to_json()/from_json()} quản lý metadata
sidecar.

Các helper này chỉ enforce expiry nếu đúng verifier được gọi ở mọi điểm chấp
nhận key. Demo native hiện cấp trực tiếp một raw 32-byte key và không gọi
\path{KeyCertificate.verify()}; vì vậy certificate module không tự làm cho raw
key native hết hạn.

\section{Demo đã chạy: input cụ thể và cách suy luận}

Session kiểm chứng gần nhất nằm ở
\path{demo/runs/session_20260925T033640Z/}. Cover là
\path{data/encoded/foreman_cif_q18_g1_300f.h264}; message 15 byte; proof nén
129 byte; tất cả stage log 00--13 báo PASS trong \path{SUMMARY.txt}.

\begin{longtable}{p{.27\textwidth}p{.31\textwidth}p{.32\textwidth}}
\toprule\textbf{Artifact/log} & \textbf{Giá trị thực tế} & \textbf{Có thể suy ra / chưa thể suy ra} \\
\midrule\endfirsthead
\toprule\textbf{Artifact/log} & \textbf{Giá trị thực tế} & \textbf{Có thể suy ra / chưa thể suy ra} \\
\midrule\endhead
\path{01_generate_proof.log} & 129 B proof, 513 public signal scalars & Proving chạy với setup hiện có; không chứng minh key/proving key được quản lý production. \\
\path{02_pack_payload.log} & 15 B message + 129 B proof + 4 B length = 148 B & Outer payload format đúng, 1.184 bit trước native frame overhead. \\
\path{04_candidate_capacity.log} & 300 IDR; 105,516 raw sign candidates; capacity bị cap 19,200 bit & Fixture có nhiều candidate hơn yêu cầu 1,336 bit; không đại diện mọi H.264/GOP. \\
\path{04_first_candidate_location.log} & NAL 3, MB 0, RBSP bit 91 & Parser tìm được một sign sample có địa chỉ; file probe chỉ minh họa một bit. \\
\path{05_embed_video.log} & 148 B payload; H.264 output 3,551,610 B & Native embed hoàn tất với fixture; kích thước output giống input trong sample. \\
\path{06_ffmpeg_decode.log} & FFmpeg -xerror exit 0 & Stream giải mã được; không suy ra PSNR/SSIM hay invisibility. \\
\path{07_blind_extract.log} & Trích được 148 B, không cấp cover file & Blind extraction hoạt động cho fixture/key này. \\
\path{09_verify_proof.log} & Original message valid; message sửa một byte invalid & Proof verification và negative case chạy; không kiểm mọi loại statement/tamper. \\
\path{10_wrong_key_rejection.log} & Exit 2, authenticated frame version invalid & Sai key bị reject; lỗi xảy ra ở native frame/extraction chứ không phải Groth16. \\
\path{11_http_transport.log} & HTTP embed/extract round-trip; verify endpoint 501 đúng thiết kế test & Channel vận chuyển payload hoạt động; native server chưa verify Groth16. \\
\path{12_websocket_e2e.log} & 5/5 native channel E2E tests PASS & Loopback tests xác nhận protocol/process path, không phải thiết bị camera edge production. \\
\path{13_python_pipeline.log} & Nhánh Python verify PASS, 1,176 bit; output 16,768,616 B trên deadline 600-frame input & Chứng minh nhánh legacy hoạt động với benchmark contract khác; không phải so sánh cùng nguồn hay performance công bằng. \\
\bottomrule
\end{longtable}

\subsection{Artifacts nên giữ khi debug}

\begin{itemize}
  \item \path{session.json}, \path{00_prepare.log}: xác nhận input đúng và
  fingerprint có thể tái lập.
  \item \path{candidate_capacity.json}, \path{04_positions_summary.txt}:
  kiểm tra payload có vượt sức chứa hay không.
  \item \path{stego_video.h264}: output chính của đường native; so hash với
  \path{http_stego_video.h264} khi kiểm tra HTTP artifact.
  \item \path{extracted_payload.bin}, \path{recovered_message.bin},
  \path{recovered_proof_129_bytes.bin}: xác định lỗi nằm ở extraction, framing,
  unpack hay ZK verify.
  \item \path{SUMMARY.txt}: trạng thái tổng hợp. Luôn mở log gốc để xem vì sao
  step fail; summary chỉ phân loại theo marker PASS/exit code.
\end{itemize}

\path{secret_key.bin} là dữ liệu nhạy cảm, dù là key demo. Không commit, gửi
người khác hoặc đưa vào log. \path{wrong_key.bin} chỉ phục vụ negative test.

\section{Các giới hạn cần nhớ khi đọc code}

\begin{enumerate}
  \item \textbf{Không phải mọi H.264 đều được hỗ trợ như nhau.} Native path nhắm
  Baseline/CAVLC và các điều kiện parser/slice cụ thể; CABAC, FMO, MBAFF và nhiều
  trường hợp profile khác nằm ngoài primitive đã mô tả.
  \item \textbf{I/P/B và nhúng là hai lớp khác nhau.} Parser native thu candidate
  từ IDR I-slice CAVLC; demo fixture all-intra khiến mỗi frame phù hợp là IDR.
  Đừng suy ra từ đó rằng mọi P/B residual hiện đều có thể làm carrier.
  \item \textbf{Blind không có nghĩa không cần key.} Blind có nghĩa extractor
  không cần cover gốc; vẫn cần key để đồng bộ lịch và authenticate frame.
  \item \textbf{Groth16 và HMAC khác nhiệm vụ.} Proof kiểm relation circuit;
  HMAC kiểm frame payload theo shared key; bearer token xác thực API caller; các
  lớp không thay thế nhau.
  \item \textbf{Native streaming và camera realtime không đồng nghĩa.} WebSocket
  nhận byte stream Annex-B và có bounded I/O; camera capture/device driver và
  chứng minh deadline trên edge cần test riêng.
  \item \textbf{Nhánh Python sinh sidecar.} \path{.positions.json},
  \path{.meta.json}, \path{.manifest.json} là artifact của nhánh cũ, không cần
  cho native blind extraction.
\end{enumerate}

\section{Lộ trình đọc source cho người mới tiếp nhận}

Để nắm hệ thống mà ít bị phân tán, nên đọc theo thứ tự sau:
\begin{enumerate}
  \item Chạy/đọc \path{demo/README.md}, rồi theo các script 00--09. Mở mỗi log
  ngay sau khi chạy để liên hệ input/output.
  \item Đọc \path{demo/05_embed_video.py} và \path{demo/07_blind_extract.py}
  để thấy Python chỉ làm orchestration ở luồng chính.
  \item Đọc \path{src/api/native_handlers.py::embed_handler()},
  \path{extract_handler()} và \path{native_stream()} để nắm boundary HTTP/WS.
  \item Đọc CLI \path{native/tools/blind_bits.cpp::main()}, sau đó
  \path{run_live_authenticated_embed()} / \path{run_live_authenticated_extract()}.
  \item Đọc C++ library theo chuỗi
  \path{split_annex_b()} $\to$ SPS/PPS/slice parsing $\to$
  \path{decode_baseline_i_idr_slices()} $\to$
  \path{collect_cavlc_trailing_one_sign_candidates()} $\to$
  \path{select_keyed_cavlc_sign_candidates()} $\to$
  \path{AuthenticatedCavlcStreamEncoder/Decoder}.
  \item Đọc proof path: \path{src/zk_proof.py::pack()},
  \path{proof_to_bytes()}, \path{ZKSnarkBridge} và
  \path{circuits/payload_verify.circom}.
  \item Cuối cùng đọc \path{src/embedder.py}/\path{src/verifier.py} để hiểu vì
  sao nhánh Python khác native; đọc \path{demo/13_python_pipeline.py} như một
  integration test/benchmark path riêng.
\end{enumerate}

\section{Tra cứu nhanh theo câu hỏi debug}

\begin{longtable}{p{.37\textwidth}p{.54\textwidth}}
\toprule\textbf{Triệu chứng} & \textbf{Nơi bắt đầu lần theo} \\
\midrule\endfirsthead
\toprule\textbf{Triệu chứng} & \textbf{Nơi bắt đầu lần theo} \\
\midrule\endhead
Không thấy video NAL/IDR & \path{demo/03_read_video.log}, C++ \path{split_annex_b()}, \path{decode_baseline_i_idr_slices()}, SPS/PPS parser. \\
Capacity thấp hoặc payload không fit & \path{04_candidate_capacity.log}, \path{collect_cavlc_trailing_one_sign_candidates()}, max bits mỗi IDR. \\
Embed trả nonzero hoặc output thiếu & \path{05_embed_video.log}, \path{run_live_authenticated_embed()}, encoder \path{process_segment()/complete()}. \\
FFmpeg từ chối video & \path{06_ffmpeg_decode.log}; xác định parser/patch/EBSP boundary trước khi kiểm quality. \\
Extract báo sai version/tag & So key bytes hai phía, cùng native schedule/frame domain, sau đó \path{AuthenticatedCavlcStreamDecoder::consume_segment()}. \\
Payload extract được nhưng ZK false & So \path{packed_payload.bin} với output unpack; kiểm 4-byte length, proof 129 B, message/key/public signals và verification key. \\
HTTP status 501 cho verify & Đúng với native adapter hiện tại: \path{create_native_app(... verify_handler=None)}. ZK verify phải chạy riêng. \\
Python chạy lâu/RAM cao & Nhánh \path{src/embedder.py}, parser toàn-video, analysis cache và \path{demo/13_python_pipeline.py}; không quy kết cho native CLI. \\
Khóa hết hạn vẫn dùng được & Kiểm tra có gọi \path{KeyCertificate.verify()} tại enforcement boundary hay chỉ cấp raw key cho native. Hiện demo path cấp raw key. \\
\bottomrule
\end{longtable}

\section{Kết luận vận hành}

Khi đọc một đường lỗi, lần theo artifact tuần tự: đầu vào/session fingerprint,
proof bytes, packed blob, capacity, stego output, extracted blob, unpacked
message/proof, rồi Groth16 result. Điểm nối quan trọng nhất là native CLI nhận
payload bytes opaque; Python demo tạo proof và blob, còn H.264 syntax parsing,
keyed sign schedule, HMAC frame, patch, blind extraction thuộc C++.

Đường native đã được chạy end-to-end trên fixture hỗ trợ, nhưng source vẫn có
nhánh Python legacy, scheduler abstraction, certificate/manifest helpers và
trust experiments không tự động nằm trong đường runtime đó. Khi thay đổi code,
hãy kiểm tra đúng consumer kế tiếp của hàm vừa sửa và chạy test ở boundary tương
ứng; PASS ở một nhánh không xác nhận các nhánh còn lại.

\end{document}
